\subsubsection{面向小 Batch 场景的量化路径优化}

为了重新验证量化路径在小 batch 区间的真实收益，本文针对 `input=4096, output=256` 的长上下文生成设置，使用隔离式 batch-size runner 对五组配置进行了重跑：对每个 `BS` 单独加载模型，并以 `3` 次重复测量吞吐量。量化路径仍沿用当前的定向运行时调优策略，即在两条 2bit 路径上启用 `quant_k_use_meta=true`，并按 batch size 对 `quant_v_split_k` 与 `quant_v_tile_config` 做运行时覆写；其中 `Sparse50+2bit` 在 `BS=3,4,6,7` 使用 `split_k=8 + tile64`，`Sparse70+2bit` 在 `BS=3,4,6,7,8` 使用 `split_k=8 + tile64`、在 `BS=5` 使用 `split_k=4 + tile64`。实验中仍未纳入 `quant_v_decode_n1=True`，因为此前验证表明其在当前设置下带来稳定负收益。

实验结果如表~\ref{tab:smallbs-quant-opt-throughput} 与图~\ref{fig:smallbs-quant-opt-throughput} 所示。与早先归档结果相比，本轮重跑采用了更严格的测量口径，因此数值趋势也更加平滑。整体来看，量化路径并没有在 `BS=1..8` 的全区间持续统治，而是呈现出“在小到中等 batch 区间最有优势、在更大 batch 下被 sparse-only 重新追平甚至反超”的特征。

从左图可以看到，在 `BS=1` 时 Dense 仍然是最优配置，吞吐量为 `28.55` tokens/s；但从 `BS=2` 开始，`Sparse50+2bit` 已明显超过 Dense 和 `Sparse50`，在 `BS=2` 时达到 `42.76` tokens/s，在 `BS=4` 时达到 `64.70` tokens/s。与此同时，`Sparse50+2bit` 相对 `Sparse50` 的优势主要集中在 `BS=1..5`：例如在 `BS=5` 时前者为 `72.02` tokens/s，而后者为 `66.93` tokens/s。然而随着 batch size 继续增大，这一优势开始收敛；在 `BS=6` 时两者已经基本持平，到了 `BS=7` 和 `BS=8`，`Sparse50` 分别以 `86.52` 和 `94.51` tokens/s 重新超过 `Sparse50+2bit` 的 `80.47` 和 `86.98` tokens/s。

从右图可以看到，`Sparse70+2bit` 的表现比 `Sparse50+2bit` 更稳定。在 `BS=2` 时其吞吐量为 `42.37` tokens/s，已高于 Dense 的 `39.81` tokens/s 和 `Sparse70` 的 `30.56` tokens/s；在 `BS=3` 到 `BS=6` 的区间内，`Sparse70+2bit` 进一步成为整体最优配置，吞吐量分别达到 `57.81`、`66.42`、`74.22` 和 `80.38` tokens/s。即便在 `BS=7` 时，`Sparse70+2bit` 仍略高于 `Sparse70`。不过在 `BS=8` 时，`Sparse70` 以 `96.43` tokens/s 重新超过 `Sparse70+2bit` 的 `87.12` tokens/s，说明在更大 batch 条件下，sparse-only 路径的吞吐扩展能力开始抵消量化带来的额外收益。

因此，这组重跑结果表明，当前的小 batch 定向优化确实显著提升了量化路径在 `BS=2..6` 区间内的竞争力，但其收益并非在 `BS=1..8` 的全区间单调扩大。更准确的结论是：量化路径在小到中等 batch size 上更容易释放 KV Cache 压缩带来的访存优势，而在更大 batch 下，quant/dequant、元数据访问以及额外布局管理带来的开销也会被放大，从而使 sparse-only 路径重新取得优势。换言之，`quant_k_use_meta=true` 与 runtime split-k/tile 调优已经足以将量化方法推到“`BS=2..6` 局部最优”的状态，但尚不足以保证其在更大 batch 区间持续优于高效的 sparse-only 实现。

需要说明的是，本节结果仍然是在共享状态下的 `NVIDIA A100 80GB PCIe` 上获得，因此更适合作为论文草稿中的预实验材料；正式定稿时，仍建议在独占 GPU 条件下按相同配置复测并替换最终数值。

\begin{table}[t]
    \centering
    \caption{Llama-3 8B 在不同 batch size 下的五线吞吐量对比（量化路径采用小 batch 定向优化，input=4096, output=256）}
    \label{tab:smallbs-quant-opt-throughput}
    \begin{tabular}{cccccc}
    \toprule
    Batch Size & Dense & Sparse50 & Sparse70 & Sparse50+2bit & Sparse70+2bit \\
    \midrule
    1 & 28.55 & 17.16 & 16.84 & 23.24 & 23.38 \\
    2 & 39.81 & 31.75 & 30.56 & 42.76 & 42.37 \\
    3 & 49.15 & 45.15 & 46.18 & 55.45 & 57.81 \\
    4 & 55.83 & 53.80 & 57.63 & 64.70 & 66.42 \\
    5 & 56.67 & 66.93 & 65.96 & 72.02 & 74.22 \\
    6 & 56.81 & 77.17 & 78.01 & 76.80 & 80.38 \\
    7 & 59.69 & 86.52 & 83.53 & 80.47 & 84.06 \\
    8 & 62.80 & 94.51 & 96.43 & 86.98 & 87.12 \\
    \bottomrule
\end{tabular}

\end{table}

\begin{table}[t]
    \centering
    \caption{优化后量化路径在小 batch 区间的代表性表现}
    \label{tab:smallbs-quant-opt-best}
    \begin{tabular}{ccccc}
        \toprule
        Batch Size & 最优配置 & 吞吐量 & 备注 & 说明 \\
        \midrule
        1 & Dense & 28.55 & 小 batch 极限点 & Dense 仍保持最高吞吐量 \\
        2 & Sparse50+2bit & 42.76 & 量化反超起点 & 量化路径首次成为整体最优配置 \\
        4 & Sparse70+2bit & 66.42 & 关键优化点 & 量化在小到中等 batch 区间收益最明显 \\
        6 & Sparse70+2bit & 80.38 & 量化领先末端 & 这是量化仍保持全局最优的最后一个 batch 点 \\
        8 & Sparse70 & 96.43 & 大 batch 回切点 & sparse-only 在更大 batch 下重新取得最优 \\
        \bottomrule
    \end{tabular}
\end{table}

\begin{figure}[t]
    \centering
    \begin{minipage}[t]{0.49\linewidth}
        \centering
        \includegraphics[width=\linewidth]{figures/throughput_vs_bs_dense_sparse50_quant50.pdf}
        
        \vspace{2mm}
        \small (a) Dense / Sparse50 / Sparse50+2bit
    \end{minipage}\hfill
    \begin{minipage}[t]{0.49\linewidth}
        \centering
        \includegraphics[width=\linewidth]{figures/throughput_vs_bs_dense_sparse70_quant70.pdf}
        
        \vspace{2mm}
        \small (b) Dense / Sparse70 / Sparse70+2bit
    \end{minipage}
    \caption{Llama-3 8B 在不同 batch size 下的三线吞吐量对比（量化路径采用小 batch 定向优化，input=4096, output=256）。左图展示 Dense、Sparse50 与 Sparse50+2bit；右图展示 Dense、Sparse70 与 Sparse70+2bit。}
    \label{fig:smallbs-quant-opt-throughput}
\end{figure}

综合而言，本节结果表明，经过针对量化 `Key` 与 `Value` 路径的联合调优后，2bit 量化方法已经从“仅在局部 batch size 上具备竞争力”提升为“在 `BS=2..6` 区间具备明确优势”的状态。其中，`BS=1` 时 Dense 仍然保持最高吞吐量；从 `BS=2` 开始，量化路径先后在 `BS=2` 由 `Sparse50+2bit`、在 `BS=3..6` 由 `Sparse70+2bit` 取得全局最优；而在 `BS=7` 和 `BS=8` 时，最优配置又分别回到 `Sparse50` 与 `Sparse70`。这说明在当前 `4096 \rightarrow 256` 的设置下，量化 KV Cache 的端到端收益主要集中在小到中等 batch 区间，后续若希望在更大 batch 下继续保持领先，仍需要进一步压缩量化路径自身的附加开销。

\subsubsection{压缩模块时间与压缩率}

为了在与端到端实验相同的工作负载下评估压缩模块本身的行为，本文进一步在 `Meta-Llama-3-8B-Instruct` 上对 prompt-side KV Cache 压缩过程进行了统计。实验设置保持与前文一致，即 `input=4096, output=256, BS=1..8`，并在 `Sparse50`、`Sparse70`、`Sparse50+2bit` 与 `Sparse70+2bit` 四条路径上记录每次 prefill 后的原始 KV 大小、压缩后 KV 大小以及压缩阶段的总耗时。由于该统计直接插入到真实模型的 prefill 压缩路径中，因此其结果能够反映压缩模块在端到端工作负载下的实际开销，而不是脱离推理场景的独立 microbenchmark。

实验结果如图~\ref{fig:smallbs-compression-stats} 和表~\ref{tab:smallbs-compression-summary} 所示。首先，从压缩率来看，量化路径相较 sparse-only 路径具有显著优势，而且这一优势在不同 batch size 下基本稳定：`Sparse50` 的压缩率约为 `1.51\times`，`Sparse70` 为 `2.09\times`；引入 2bit 量化后，`Sparse50+2bit` 提升到 `3.25\times`，`Sparse70+2bit` 进一步提升到 `3.48\times`。以 `BS=8` 为例，原始 KV Cache 总大小约为 `4096` MB，`Sparse50` 可压缩到 `2720.16` MB，而 `Sparse50+2bit` 可进一步压缩到 `1261.16` MB；对应地，`Sparse70` 与 `Sparse70+2bit` 则分别压缩到 `1964.29` MB 与 `1176.25` MB。也就是说，在真实长上下文设置下，2bit 量化相较 sparse-only 版本仍然带来了额外的 KV Cache 存储压缩收益。

其次，从压缩模块时间来看，四条路径的压缩时间都随 batch size 增长近似线性上升，但量化版本在同一 sparsity 下始终更快。对于 50\% 稀疏路径，`Sparse50` 的压缩时间从 `BS=1` 的 `740.80` ms 增长到 `BS=8` 的 `5688.23` ms，而 `Sparse50+2bit` 则从 `545.18` ms 增长到 `4317.88` ms；对于 70\% 稀疏路径，`Sparse70` 的压缩时间从 `716.30` ms 增长到 `5411.29` ms，而 `Sparse70+2bit` 则从 `552.86` ms 增长到 `4293.18` ms。可见，量化路径不仅压缩率更高，而且在 prompt-side 压缩阶段本身也没有引入额外的大额时间惩罚，反而由于压缩后布局更紧凑，整体压缩耗时略低于 sparse-only 版本。

综合来看，这组结果说明：量化路径在当前工作负载下同时提供了更高的 KV Cache 压缩率与更低的 prompt-side 压缩时间。其中，`Sparse70+2bit` 在四条压缩路径中始终取得最高压缩率，而两条 2bit 路径在压缩时间上的优势则说明，量化方法的端到端收益并不仅来自后续 decode 阶段的访存节省，也来自压缩模块本身更高效的表示与写回方式。

\begin{table}[t]
    \centering
    \caption{代表性 batch size 下的 KV Cache 压缩率与压缩时间（input=4096）}
    \label{tab:smallbs-compression-summary}
    \begin{tabular}{lcccccc}
    \toprule
    Config & BS1 Ratio & BS1 Time & BS4 Ratio & BS4 Time & BS8 Ratio & BS8 Time \\
    \midrule
    Sparse50 & 1.51x & 741 & 1.51x & 3095 & 1.51x & 5688 \\
    Sparse70 & 2.09x & 716 & 2.09x & 2890 & 2.09x & 5411 \\
    Sparse50+2bit & 3.25x & 545 & 3.25x & 2122 & 3.25x & 4318 \\
    Sparse70+2bit & 3.48x & 553 & 3.48x & 2126 & 3.48x & 4293 \\
    \bottomrule
\end{tabular}

\end{table}

\begin{figure}[t]
    \centering
    \begin{minipage}[t]{0.49\linewidth}
        \centering
        \includegraphics[width=\linewidth]{figures/compression_time_vs_bs.pdf}
        
        \vspace{2mm}
        \small (a) Prompt-side compression time
    \end{minipage}\hfill
    \begin{minipage}[t]{0.49\linewidth}
        \centering
        \includegraphics[width=\linewidth]{figures/compression_ratio_vs_bs.pdf}
        
        \vspace{2mm}
        \small (b) KV Cache compression ratio
    \end{minipage}
    \caption{小 batch 工作负载下压缩模块的时间与压缩率统计。左图展示 prompt-side 压缩时间随 batch size 的变化；右图展示压缩率（原始 KV 大小 / 压缩后 KV 大小）随 batch size 的变化。}
    \label{fig:smallbs-compression-stats}
\end{figure}
